\begin{table}[H]
\centering
\caption*{\textbf{Quadro 2 - Posição da MOSAIC em relação aos trabalhos-base}}
\footnotesize
\renewcommand{\arraystretch}{1.16}
\begin{tabularx}{\linewidth}{|p{0.15\linewidth}|p{0.18\linewidth}|p{0.27\linewidth}|Y|}
\hline
\textbf{Trabalho} & \textbf{Unidade principal} &
\textbf{Relação entre tarefa e skill} & \textbf{Decisão aproveitada pela MOSAIC} \\
\hline
Agent Skills & arquivo ou diretório de instruções & ativação de uma skill &
\texttt{SKILL.md} como fonte canônica \\
\hline
SkillRouter & consulta de retrieval & seleção e reranking de skills &
body completo como sinal de roteamento \\
\hline
SkillWeaver & subtarefa atômica & uma skill por subtarefa &
feedback de retrieval para a decomposição \\
\hline
ReAct & trajetória de ação e observação & não define catálogo de skills &
revisão após observações \\
\hline
MOSAIC 0.2 & objetivo orientado a resultado & zero, uma ou várias skills por objetivo &
composição many-to-many e evidência capturada pelo runtime \\
\hline
\end{tabularx}
\par\smallskip\raggedright\small Fonte: elaboração própria (2026).
\end{table}
